\documentclass[11pt,a4paper]{article}
\usepackage[margin=1.05in]{geometry}
\usepackage{amsmath,amssymb,amsthm}
\usepackage{microtype}
\usepackage[hidelinks]{hyperref}
\hypersetup{pdftitle={From the 3-free door to a Hecke $L$-function},pdfauthor={Jack Pickett}}

\newtheorem{lemma}{Lemma}
\newtheorem{theorem}{Theorem}
\newtheorem{corollary}{Corollary}
\newtheorem{remark}{Remark}

\title{From the 3-free door to a Hecke $L$-function}
\author{Jack Pickett \textperiodcentered\ 9th October 2026}
\date{}

\begin{document}
\maketitle

\begin{abstract}
The 3-free door \(m_0(p)=p+\chi_3(p)\) of an odd prime \(p\neq 3\) selects one of the two even neighbours of \(p\). The signed imbalance between the two doors is measured by a Dirichlet character \(C=\chi_3\psi_5\). Absolute norm lifts \(C\) to a multiplicative function \(\eta\) on the nonzero ideals of the Eisenstein integers \(\mathcal{O}_K\) for \(K=\mathbb{Q}(\sqrt{-3})\). The function takes values in the fourth roots of unity, is trivial on the principal ray class \(1+15\mathcal{O}_K\), and therefore descends to a finite-order character of the ray class group modulo \(15\mathcal{O}_K\). The associated ideal series converges absolutely for \(\operatorname{Re}(s)>1\) and admits the factorization
\[
L_K(s,\eta)=L(s,C)\,L(s,\mathrm{companion15})
\]
as entire functions, where the companion differs from the pure quartic character \(\psi_5\) by the explicit Euler factor \(1+i\,3^{-s}\). In the half-plane \(\operatorname{Re}(s)>0\) the zeros of \(L_K(s,\eta)\) are precisely the zeros of either factor. The resulting Hecke \(L\)-function is admissible for every fixed zero-free half-plane known for finite-order characters of \(K\).
\end{abstract}

\section{The door character}

Let \(\chi_3\) be the non-principal Dirichlet character modulo \(3\), equivalently the Kronecker symbol \((-3/\cdot)\). For an odd prime \(p\neq 3\) the two even neighbours are
\[
m_0(p)=p+\chi_3(p),\qquad m_1(p)=p-\chi_3(p).
\]
Exactly one is divisible by \(3\); the other is the 3-free door. The product identity \(m_0(p)m_1(p)=p^2-1\) holds, and both doors are even. The signed imbalance between the doors is recorded by the Dirichlet character
\[
C=\chi_3\psi_5,
\]
where \(\psi_5\) is the primitive quartic character modulo \(5\). Twisting by \(\chi_3\) produces the companion character, which agrees with \(\psi_5\) away from multiples of \(3\). The two \(L\)-functions therefore differ by the explicit Euler factor at the ramified prime:
\begin{equation}
\label{eq:correction}
L(s,\mathrm{companion15})=(1+i\,3^{-s})L(s,\psi_5).
\end{equation}
The correction is holomorphic and non-zero for \(\operatorname{Re}(s)>0\), so the companion and the pure quartic share their zeros in that half-plane.

\section{The ideal character}

Let \(K=\mathbb{Q}(\sqrt{-3})\) and let \(\mathcal{O}_K\) be its ring of integers, the Eisenstein integers. Absolute norm defines a map from nonzero ideals of \(\mathcal{O}_K\) to positive integers. Composing with \(C\) yields a multiplicative function
\[
\eta(I)=C(N(I))
\]
on the nonzero integral ideals, with the value zero when the ideal meets the modulus. The same formula extends to fractional ideals coprime to \(15\mathcal{O}_K\).

\begin{lemma}
For every ideal \(I\) coprime to \(15\mathcal{O}_K\),
\[
\eta(I)^4=1.
\]
In particular the values are nonzero fourth roots of unity.
\end{lemma}

\begin{lemma}
If \(a\in\mathcal{O}_K\) satisfies \(a\equiv 1\pmod{15\mathcal{O}_K}\), then \(\eta((a))=1\).
\end{lemma}

The two statements supply the data of a finite-order character of the ray class group modulo \(15\mathcal{O}_K\). Extending the values to fractional ideals and verifying invariance under the full ray-class equivalence relation produces a homomorphism from the group of fractional ideals coprime to the modulus onto the fourth roots of unity, trivial on the principal ray class. The resulting character is denoted \(\eta\) below.

\section{The ideal series}

The Hecke series attached to \(\eta\) is
\[
L_K(s,\eta)=\sum_{I\neq 0}\frac{\eta(I)}{N(I)^s},
\]
summed over the nonzero integral ideals of \(\mathcal{O}_K\). Absolute convergence for \(\operatorname{Re}(s)>1\) follows from the bound \(|\eta(I)|\le 1\) and the linear growth of the ideal-counting function. The number of ideals of norm \(n\) equals the divisor sum
\begin{equation}
\label{eq:coeff}
\#\{I:N(I)=n\}=\sum_{d\mid n}\chi_3(d).
\end{equation}
The identity is obtained by combining the exact prime-power counts (split, inert, and ramified) with coprime multiplicativity of the counting function.

\section{Factorization and continuation}

Comparing Euler factors at split, inert, and ramified primes, and using the coefficient identity \eqref{eq:coeff}, yields the factorization
\begin{equation}
\label{eq:factor}
L_K(s,\eta)=L(s,C)\,L(s,\mathrm{companion15})
\end{equation}
for \(\operatorname{Re}(s)>1\). The correction \eqref{eq:correction} accounts for the distinction at the ramified prime \(3\). Both sides of \eqref{eq:factor} admit entire continuations, and the continuation is unique. Consequently the identity holds as an equality of entire functions.

\begin{theorem}
In the half-plane \(\operatorname{Re}(s)>0\) the zeros of \(L_K(s,\eta)\) are precisely the zeros of \(L(s,C)\) or of \(L(s,\psi_5)\). The function \(L_K(s,\eta)\) is non-vanishing for \(\operatorname{Re}(s)\ge 1\), and every point on the line \(\operatorname{Re}(s)=1\) has a zero-free neighbourhood.
\end{theorem}

The non-vanishing on \(\operatorname{Re}(s)\ge 1\) is inherited from the classical non-vanishing of the Dirichlet factors; the local disks on the boundary follow by continuity.

\section{Admissibility}

The character \(\eta\) is a finite-order Hecke character of \(K\). Every fixed zero-free half-plane established for the family of finite-order Hecke \(L\)-functions over \(K\) therefore applies to \(L_K(s,\eta)\). In particular the half-plane \(\operatorname{Re}(s)>7/8\) and the half-plane \(\operatorname{Re}(s)>47/48\) both exclude zeros of \(L_K(s,\eta)\), and hence exclude zeros of both Dirichlet factors in those regions. The construction supplies an explicit character, originating from the 3-free door, that is admissible for these theorems. No improvement of the exponents is claimed.

\begin{thebibliography}{9}
\bibitem{openai}
OpenAI, A zero-free half-plane \(\operatorname{Re}s>7/8\), 30 September 2026.
\bibitem{kintali}
S.~Kintali, A short proof of the Quasi-Riemann Hypothesis, 7 October 2026.
\end{thebibliography}

\end{document}
